% ==========================================
% CHAPTER 9: GLOSSARY
% ==========================================
\chapter{Glossary} % 

\begin{description}
    \item[Argument] An argument is essentially an input to a function. %  
    Within python, functions, both built in and user defined, are referenced much like in the math example: via parenthesis attached to the function call. %  

    \item[CLI] Command-line Interface (terminal). A means of interacting with a computer system via successive lines of text in the form of commands. % 

    \item[Command] A word, phrase, or instruction that can be understood and interpreted by a computer system which then executes the command in question. % 

    \item[Conditional Statement] A statement defining a certain condition, using operators like greater than, less than, equals to, their opposites, or some combination. %  
    These statements enclose a block of code that runs only if the conditional statement is evaluated to be true. % 

    \item[Data Type] Refers to the different types of objects in python to which python places certain rules. %  
    For example, integers and floats can be added, but not indexed. % 

    \item[Directory] Within a UNIX system, a directory is analogous to ``folders'' on a PC or Mac. %  
    It's where all your files are stored. % 

    \item[Documentation] A special string added when defining functions which specifies the arguments of the function and their data types, accessible via \texttt{help(function\_name)}. % 

    \item[Element] A single discrete object within an iterable set. For example, a single character in a string, or a single entry in a list or array, or an entry in a dictionary. % 

    \item[FITS] Flexible Image Transport System. A file format typical to astronomical images. % 

    \item[Flag] Also known as an option, a flag is a way to modify a UNIX command to alter the way it performs a task. %  
    The typical syntax for a flag is a dash (-) followed by a letter or short combination thereof. % 

    \item[For-loop] A for loop is a block of code that contains some iterative variable like ``i'' within it, with ``i'' cycling through different values defined in the initiation of the loop. % 

    \item[Function] A function is an operator that takes some inputs (or none) and, when called, performs some operations and outputs something (or multiple things). % 

    \item[GUI] Graphical User Interface. A means of interacting with a computer system via graphical icons and visual indicators, through the direct manipulation of graphical elements. % 

    \item[Index] Given a list, array, or other sliceable data type, every element is assigned an index based on its position from the leftmost element, starting with 0. % 

    \item[Library] A large collection of functions that can be used in Python by importing the library. %  
    Calling functions from libraries usually requires the dot notation call of the library name (dot) function name. % 

    \item[Loop] A block of code that is run multiple times, either due to iteration through some predefined range of values, or indefinitely so long as some condition is met. % 

    \item[Operating System] Software that manages the hardware and software resources for a computer. %  
    Common OS's include Windows, MacOS, and Ubuntu. % 

    \item[Root Directory] The root directory is the directory within which all other folders/directories and files are contained. % 

    \item[Syntax] The specific set of words, phrases, and commands that can be successfully interpreted and understood by a computer. % 
\end{description}